\subsection{Warstwa 1: testy programistyczne (Unit + Integration)}

Testy warstwy programistycznej zostaly wykonane w srodowisku lokalnym projektu z wykorzystaniem interpretera \texttt{C:\textbackslash Development\textbackslash ContentSaaS\textbackslash backend\textbackslash .venv\textbackslash Scripts\textbackslash python.exe}. Do uruchomienia pakietu testow zastosowano narzedzie \texttt{pytest} z pluginem \texttt{pytest-cov}.

Bazowa komenda pomiarowa miala postac: \texttt{python -m pytest backend/tests --cov=backend --cov-report=html --cov-report=term}. Zakres ten obejmuje testy jednostkowe oraz integracyjne backendu, w tym testy API, testy warstwy przetwarzania tresci, testy warstwy persystencji oraz testy narzedzi pomocniczych.

W uruchomieniu zebrano 53 testy i wszystkie zakonczyly sie statusem PASS (0 FAIL, 0 ERROR). Otrzymany wynik potwierdza stabilnosc biezacej implementacji w zakresie pokrytych scenariuszy funkcjonalnych i walidacyjnych.

\begin{table}[htbp]
\centering
\caption{Wynik uruchomienia testow Unit + Integration}
\label{tab:unit_integration_summary}
\begin{tabular}{lr}
\hline
Metryka & Wartosc \\
\hline
Liczba testow & 53 \\
PASS & 53 \\
FAIL & 0 \\
ERROR & 0 \\
Coverage surowe (\texttt{--cov=backend}) & 58\% \\
Coverage modulow produkcyjnych & 50\% \\
\hline
\end{tabular}
\end{table}

Wynik coverage wymaga rozroznienia dwoch perspektyw metrycznych. Surowe 58\% obejmuje wszystkie pliki pod \texttt{backend/}, lacznie z plikami testowymi. Dla oceny jakosci testowania kodu produkcyjnego zastosowano dodatkowo raport filtrowany (\texttt{--include="backend/*.py"} oraz \texttt{--omit="backend/tests/*,backend/test\_*.py"}), ktory daje wynik 50\%.

W dalszej tresci pracy jako glowny wskaznik przyjeto 50\%\footnote{Wartosc obejmuje wylacznie moduly produkcyjne backendu. Wynik 58\% jest metryka surowa i zawiera pliki testowe.}, poniewaz ten licznik lepiej reprezentuje realne pokrycie kodu dzialajacego w systemie. Roznica miedzy 58\% i 50\% wynika z wysokiego pokrycia samych plikow testowych.

Wyniki raportu coverage zostaly udokumentowane graficznie. Rysunek~\ref{fig:coverage_report_summary} przedstawia widok zbiorczy raportu, a Rysunek~\ref{fig:coverage_report_full} prezentuje pelna liste plikow z poziomem pokrycia.

\begin{figure}[htbp]
\centering
\includegraphics[width=0.95\linewidth]{coverage-report.png}
\caption{Raport coverage: widok zbiorczy.}
\label{fig:coverage_report_summary}
\end{figure}

\begin{figure}[htbp]
\centering
\includegraphics[width=0.95\linewidth]{coverage-report-full.png}
\caption{Raport coverage: pelna lista plikow.}
\label{fig:coverage_report_full}
\end{figure}

Analiza wynikow wskazuje, ze najwyzsze pokrycie uzyskaly moduly rdzeniowe przetwarzania i persystencji (m.in. \texttt{content\_processor.py} i \texttt{database.py}), natomiast najnizsze wartosci dotycza plikow narzedziowych, debugowych i benchmarkowych oraz czesci scenariuszy I/O i integracji z zasobami zewnetrznymi.
